% Tapis de Quillen, with Quillen's offprint and preprint (cote 162-5, typed part): archive cote 162-5, site page 38 (PDF page 39).
% typescript. Draft compared: none.
% AI-transcribed by Claude Opus 5.5 (claude-opus-5-5, adaptive thinking, effort high), Message Batches API, batch t162_full (msgbatch_01UAh3zyWPd454AeDkhnmQ33), 2026-11-13; not yet checked by eye.
\scanpage{162-5/38}
\margin{! Attention à la catégorie \uline{Hom} grab, qui s'\unsure{identitifie|identifie} à la catégorie \unsure{définie} par un $R\mathrm{Hom}_{Z}(C_{*},C_{*})$, ce point devrait être élucidé!}

de même les (Cat)-groupes(ou groupes à valeurs dans (Cat)) comme étant
les groupes semisimpliciaux qui "sont"des catégories i.e. qui \xout{sexdéfinisx}
\xout{sent} satsifont la condition d'exactitude familière, permettant d'en déter
miner tous les composants à l'aide de C$_{o}$\xout{$_{x}$} et C$_{1}$ . Appliquant le tapis de
Dold, dans le cas commutatif, on trouve que les (Cat)-groupes abéliens
forment une catégorie équivalente à celle des complexes de chaines de
groupes abéliens qui sont de longueur 1:
\[ 0 \leftarrow C_{o} \leftarrow C_{1} \leftarrow 0 \quad . \]
Quand on explicite \unsure{la|le} (Cat)-groupe abélien associé à un tel complexe de
chaines, on trouve "ma" construction chérie \unsure{da|de} la "catégorie de Picard"
associée au complexe: objets ceux de C$_{o}$ , flêches de x dans y les z dans C$_{1}$
tels que dz = y-x . On trouve donc à isomorphisme \unsure{ppès|près} tous les groupes
abéliens de (Cat) ainsi. Un complément important est que toute catégorie
\inserted[typed]{\deleted{(resp. et}}
de Picard \deleted{commutative)}est équivalente (au sens technique qu'on ne précise
\inserted[typed]{\deleted{(resp. (Cat)-groupe}}
pas ici) à \xout{laxcatégorie} un vrai (Cat)-groupe \deleted{(commutatif)} . \margin{Quand on}
\inserted{\unsure{homotopie} mod ``Gros'' équivalence, ce qui correspond à travailler dans la catégorie dérivée des complexes. Seuls invariants sont $H_0$, $H_1$ (\unsure{et l'invariant} $\pi_0\,\pi_1$).}
\margin{\unsure{Quitter définir} mais cela contredit les résultats de Verdier ? \illegible{3}}
\margin{\unsure{Mais si on prend} d'une cat. de Picard commutative, elle \unsure{ne vérifie pas} \unsure{(Cat)-groupe} \illegible{2}}
On notera qu'on a une \xout{lxxexp} classification beaucoup plus restreinte
que pour les H-espaces (resp. les H-espaces homotopiquement comm.) mod
équivalence d'homotopie, de façon précise entre les groupes resp. les grou-%
pes abéliens de la catégorie homotopique. La raison en est quon est beau\xout{xx}
\margin{//} coup plus restrictif dans le contexte catégorique pour le sens \unsure{da|de} l'associa
tivité (resp. et de la commutativité), puisqu'on ne demande pas seulement \margin{//}
l'existence d'isomorphismes d'associa\xout{i}tivité\xout{,} (qui suffirait pour avoir un
\inserted[typed]{la catégorie homotopique}
groupe dans \hand{$\gamma$}(Cat)' \hand{$\simeq$} (Ssimpl)', modulo l'existence d'unités et d'inverse)
mais qu'on suppose \uline{donné\unsure{\deleted{e}}s} ces isomorphismes, de façon à satisfaire à cer-%
taines compatibilités: c'est ce qui permet, me semble-t-il, de prouver qu'à
\margin{//} \margin{(} équivalence près, on peut supposer qu'il s'agit de vrais groupes dans
(Cat), \deleted{resp. de vraie commutativité.}

\mspage{\hand{38}}
